\section{Lean 4 Proof Listings}\label{app:lean}

The complete Lean 4 formalization is available at:
\begin{center}
\url{https://doi.org/10.5281/zenodo.18140966}
\end{center}

The formalization is not a transcription of the paper text. It exposes reusable definitions for reductions, decision problems, hardness transfers, and tractable subcases, so the main complexity claims can be checked mechanically.

\subsection{What the Artifact Covers}

The appendix supports the non-physical theorem stack of the paper:
\begin{itemize}
\item core decision-problem and sufficiency definitions
\item polynomial-time reduction infrastructure
\item coNP, PP, PSPACE, and $\Sigma_2^P$ hardness reductions
\item shifted-family approximation-gap and counted-runtime threshold constructions
\item ETH-conditioned dichotomy statements
\item cross-regime bridge lemmas (static/stochastic/sequential)
\item tractable subcases under explicit structural assumptions
\end{itemize}

The point of the artifact is precision. Once the definitions are fixed, Lean verifies that the reductions preserve the intended predicates and that the later theorems invoke the correct formal statements.

\subsection{Module Structure}

The formalization consists of \LeanTotalFiles\ files organized as follows:

\begin{itemize}
\item \texttt{Basic.lean} and \texttt{Sufficiency.lean}: core definitions for decision problems, projections, optimizer maps, and sufficiency
\item \texttt{AlgorithmComplexity.lean} and \texttt{PolynomialReduction.lean}: polynomial-time reduction infrastructure and composition
\item \texttt{Reduction.lean} and \texttt{Hardness/}: the TAUTOLOGY, SET-COVER, ETH, approximation-gap, and $\Sigma_2^P$ reduction stack
\item \texttt{StochasticSequential*.lean}: PP/PSPACE regime definitions, hardness, and hierarchy bridges
\item \texttt{QueryComplexity.lean}, \texttt{Dichotomy.lean}, and \texttt{ComplexityMain.lean}: query lower bounds and summary complexity statements
\item \texttt{Tractability/}: bounded actions, separable utility, tensor-rank, tree-structured, treewidth, and symmetry arguments
\end{itemize}

\subsection{Key Theorems}

\begin{theorem}[Polynomial Composition, Lean]\label{thm:poly-compose}
Polynomial-time reductions compose to polynomial-time reductions.
\end{theorem}

\begin{verbatim}
theorem PolyReduction.comp_exists
    (f : PolyReduction A B) (g : PolyReduction B C) :
    exists h : PolyReduction A C,
      forall a, h.reduce a = g.reduce (f.reduce a)
\end{verbatim}

\begin{theorem}[TAUTOLOGY Reduction Correctness, Lean]
The canonical reduction used in Theorem~\ref{thm:sufficiency-conp} is mechanized and proves that the empty set is sufficient exactly when the source formula is a tautology.
\end{theorem}

\begin{verbatim}
theorem reduction_correct
    (φ : QBF) :
    empty_sufficient (reduce φ) ↔ QBF.eval φ
\end{verbatim}

\subsection{Verification Status}

\begin{itemize}
\item Total lines: \LeanTotalLines
\item Theorems/lemmas: \LeanTotalTheorems
\item Files: \LeanTotalFiles
\item Status: All proofs compile with \LeanTotalSorry\ \texttt{sorry}
\end{itemize}
